\scanpage{204}
Le cas $H$ quelconque s'en déduit

\deleted{Soit} Autres énoncés : $G$ simple sur $k$
$\Longleftrightarrow$ $G$ séparable sur $k$ \unsure{;|,} $\Longleftrightarrow$ $G$ séparable sur $k$ en $e$

\begin{center}$-\;-\;-\;-\;-\;-\;-$\\[-2pt]\rule{2.5cm}{0.3pt}\end{center}

Supposons $G$ \uline{commutatif}, \deleted{Soit} alors
% slanted
\margin{inutile cf plus bas}
$\Gamma=G/G^0$ est un \deleted{\unsure{préschéma}}\inserted{\uuline{discret}} $k$-groupe commutatif \inserted{séparable sur $k$}
donc $\Gamma\xrightarrow{\times n}\Gamma$ est un hom de groupe,
dont le noyau $\,{}_n\Gamma$ est un sous-groupe
fermé.

\uline{Proposition} \deleted{Il existe un plus grand entier} \inserted{Soit $\Gamma$ un $k$-\unsure{groupe|gpe}, loc. fini sur $k$. La suite des $\,{}_n\Gamma$ est}
localement constante, et \deleted{\unsure{c'est un sous}}\inserted{son sup,,}
soit $\deleted{\unsure{a}}\Gamma^\tau$, est un sous-groupe ouvert de
$\Gamma$. Pour que \deleted{$a\in$} l'élément $a$ de $\Gamma$
soit dans $\Gamma^\tau$, il faut. \unsure{so.|su.} qu'il
existe un $n$ tel que $na=e$.
% slanted
\margin{après \unsure{un|une} \unsure{suite des|mot des|mot sur les} groupes \uline{finis} et des sous-groupes $\Gamma$}

Soit $a\in\Gamma$, \deleted{tel que} si $a\in{}_n\Gamma$ on a
$na=0$. Réciproqt, supposons $na=0$,
comme il existe $m$ tel que $m\Gamma_{a\deleted{\unsure{k}}}=0$
[réf] il s'ensuit que \deleted{\unsure{mor.}} $\{a\}\deleted{\unsure{u}}\xrightarrow{\mathrm{incl.}}\Gamma\xrightarrow{mn}\Gamma$
% archive sheet number: 203
